\begin{enumerate}
		\item On considère la fonction $f$ définie sur $]0 ~;~ +\infty [$ par : $f(x)=x\, \ln(x)$.

		\textbf{Affirmation 1 :}
$\ds\int_{1}^{\e} f(x) \,\mathrm{d} x=\dfrac{\e^{2}+1}{4}$

On calcule $\ds\int_{1}^{\e} f(x) \d x = \ds\int_{1}^{\e} x\,\ln(x) \d x$ en faisant une intégration par parties:

$\ds\int_{a}^{b} u'(x)\,v(x) \d x = \left [u(x)\,v(x)\strut \right ]_{a}^{b}- \int_{a}^{b} u(x) \, v'(x) \d x$.

$\left \{
\begin{array}{r !{=} l !{\implies} r !{=} l}
u'(x) & x & u(x) & \dfrac{x^2}{2}\\[8pt]
v(x) & \ln(x) & v'(x) & \dfrac{1}{x}
\end{array}
\right .$

$\begin{aligned}
 \ds\int_{1}^{\e} x\,\ln(x) \d x 
& = \left [\dfrac{x^2}{2}\times \ln(x) \right ]_{1}^{\e}- \int_{1}^{\e} \dfrac{x^2}{2}\times \dfrac{1}{x} \d x
= \left [\dfrac{x^2}{2}\times \ln(x) \right ]_{1}^{\e}- \int_{1}^{\e} \dfrac{x}{2} \d x\\
& = \left [\dfrac{x^2}{2}\times \ln(x) \right ]_{1}^{\e}-  \left [ \dfrac{x^2}{4}  \right ]_{1}^{\e}
= \left [\dfrac{\e^2}{2}\times \ln(\e) - \dfrac{1^2}{2}\times \ln(1) \right ]-  \left [ \dfrac{\e^2}{4} - \dfrac{1^2}{4} \right ]\\
& = \dfrac{\e^2}{2} - 0 - \dfrac{\e^2}{4} + \dfrac{1}{4}
= \dfrac{\e^2}{4}+\dfrac{1}{4}
= \dfrac{\e^2+1}{4}
\end{aligned}$

\hfill\textbf{Affirmation 1 vraie}

\item Soient $n$ et $k$ deux entiers naturels non nuls tels que $k \leqslant n$.

\textbf{Affirmation 2 :}
$\ds n \times\binom{n-1}{k-1}=k \times\binom{ n}{k}$

On va utiliser les propriétés: $\dfrac{k}{k\,!} = \dfrac{1}{(k-1)\,!}$ et $n\,! = n\times (n-1)\,!$.

\begin{list}{\textbullet}{}
\item 
$\ds\binom{n}{k}= \dfrac{n\,!}{k\,! \times (n-k)\,!}$ donc

$\ds k\times \binom{n}{k}= k\times \dfrac{n\,!}{k\,! \times (n-k)\,!}
= \dfrac{k\times n\,!}{k\,! \times (n-k)\,!}
= \dfrac{n\,!}{(k-1)\,! \times (n-k)\,!}$
\item
$\begin{aligned}[t]
\ds\binom{n-1}{k-1}
& = \dfrac{(n-1)\,!}{(k-1)\,! \times \left ((n-1)-(k-1)\right )\,!}
 = \dfrac{(n-1)\,!}{(k-1)\,! \times (n-1-k+1)\,!}\\
&= \dfrac{(n-1)\,!}{(k-1)\,! \times (n-k)\,!} \text{ donc}
\end{aligned}$ 

$\ds n\times \binom{n-1}{k-1}
= n\times \dfrac{(n-1)\,!}{(k-1)\,! \times (n-k)\,!}
= \dfrac{n\times (n-1)\,!}{(k-1)\,! \times (n-k)\,!}
== \dfrac{n\,!}{(k-1)\,! \times (n-k)\,!}$
\end{list}

\hfill\textbf{Affirmation 2 vraie}

		\item Pour les trois affirmations suivantes, on considère que l'espace est muni d'un repère orthonormé \Oijk.

\medskip

Soit $d$ la droite de représentation paramétrique : \quad 
$\left\{\begin{array}{l c l}
x&=&\phantom{2}t+1 \\
y&=&2 t+1 \\
z&=&-t
\end{array},\quad t \in \R\right.$.

		Soit $d'$ la droite de représentation paramétrique : \quad $\left\{\begin{array}{l c l}
x&=&2t'-1 \\
y&=&-t'+2 \\
z&=&\phantom{-}t'+1
\end{array},\quad t' \in \R\right.$.

\medskip

Soit $P$ le plan d'équation cartésienne : $2 x+y-2 z+18=0$.

Soit A le point de coordonnées $(-1 ~;~-3 ~;~ 2)$ et B le point de coordonnées $(-5 ~;~-5~;~6)$.

On appelle plan médiateur du segment [AB] le plan passant par le milieu du segment [AB] et orthogonal à la droite (AB).

\medskip

\textbf{Affirmation 3 :} Le point A appartient à la droite $d$.

$A\in d$ si et seulement s'il existe un réel $t$ tel que:
$\left\{
\begin{array}{r !{=} l}
-1 &t+1 \\
-3 &2 t+1 \\
2 &-t
\end{array}
\right.$

Le réel $t=-2$ vérifie les trois égalités donc le point A appartient à la droite $d$.

\hfill\textbf{Affirmation 3 vraie}

\medskip

\textbf{Affirmation 4 :} Les droites $d$ et $d'$ sont sécantes.

$d$ et $d'$ sont sécantes s'il existe un couple $(t\,,\,t')$ tel que
$\left\{
\begin{array}{l !{=} r}
t+1 & 2t'-1 \\
2 t+1 & -t'+2 \\
-t & t'+1
\end{array}
\right.$

$\left\{
\begin{array}{l !{=} r}
t+1 & 2t'-1 \\
2 t+1 & -t'+2 \\
-t & t'+1
\end{array}
\right.
\iff
\left\{
\begin{array}{l !{=} r}
t & 2t'-2 \\
2 (2t'-2)+1 & -t'+2 \\
-(2t'-2) & t'+1
\end{array}
\right.$

$\iff
\left\{
\begin{array}{l !{=} r}
t & 2t'-2 \\
4t'-4+1 & -t'+2 \\
-2t'+2 & t'+1
\end{array}
\right.
\iff
\left\{
\begin{array}{l !{=} r}
t & 2t'-2 \\
5t' & 5 \\
1 & 3t'
\end{array}
\right.
\iff
\left\{
\begin{array}{l !{=} r}
t & 2t'-2 \\
t' & 1 \\
\dfrac{1}{3} & t'
\end{array}
\right.$

Le système n'a pas de solution donc les droites $d$ et $d'$ ne sont pas sécantes.

\hfill\textbf{Affirmation 4 fausse}

\medskip

		\textbf{Affirmation 5 :} Le plan $P$ est le plan médiateur du segment [AB].

Le plan médiateur du segment [AB] est l'ensemble des points M de coordonnées $(x\,,\, y \,,\, z)$ tels que $\text{MA} = \text{MB}$, ce qui équivaut à $\text{MA}^2 = \text{MB}^2$.

$\begin{aligned}
\text{MA}^2 
& = (x+1)^2+(y+3)^2 + (z-2)^2\\ 
& = x^2+2x+1 + y^2+6y+9+z^2-4z+4\\
& = x^2+2x + y^2+6y+z^2-4z+14
\end{aligned}$

$\begin{aligned}
\text{MB}^2 
& = (x+5)^2+(y+5)^2+ (z-6)^2\\
& = x^2+10x+25 + y^2+10y +25 +z^2-12z+36\\
& = x^2+10x + y^2+10y+z^2-12z+86
\end{aligned}$

$\begin{aligned}
\text{MA}^2 = \text{MB}^2
& \iff x^2+2x + y^2+6y+z^2-4z+14 = x^2+10x + y^2+10y+z^2-12z+86\\
& \iff -8x -4y + 8z -72=0
\iff -4\left ( 2x +y -2z +18\right ) = 0\\
& \iff 2x+y-2z  +18 = 0
\end{aligned}$

On retrouve l'équation du plan $P$ donc $P$ est le plan médiateur du segment [AB].

\hfill\textbf{Affirmation 5 vraie}
	\end{enumerate}
